
\section{既约概形上的\trRationalMap{}}

不严格地说，「\trRationalMap{}」就是「几乎处处有定义的\trMorphism」，正如\trRationalFunction{}%
（定义~6.6.36、练习~5.2.I）乃是「几乎处处有定义的函数」的一个名号。我们日后将会看到：
在良好的情形下，正如有理函数那样，\trRationalMap{}在其有定义之处取值是唯一的（「既约到分离」%
定理~11.3.2，Reduced-to-Separated Theorem），并且拥有一个极大的「\trDomainOfDefinition{}」（\S11.3.3）。仅就本节而言，%
\emph{我们假定 $X$ 是\trReduced{}的}，尽管如此，我们仍会尽量小心，在所有定理中把这些假设%
明说出来。一个关键的例子将是\trIrreducible\trVariety（\S7.5.7），而\trRationalMap{}的语言%
也正是用在这种场合最为频繁。

\textbf{7.5.1.}\ \emph{\trhDefinition{}。}\quad 从一个\trReduced{}\trScheme{} $X$ 到一个%
\trScheme{} $Y$ 的\textbf{\trRationalMap{}}（rational map）$\pi$，记作 $\pi\colon X\dashrightarrow Y$，
是指如下数据：从一个\trDense\trOpenSubset{} $U\subset X$ 出发的一个\trMorphism{}%
$\alpha\colon U\to Y$，并配有如下的\trEquivalenceRelation：%
$(\alpha\colon U\to Y)\sim(\beta\colon V\to Y)$，若存在一个\trDense\trOpenSet{}%
$Z\subset U\cap V$ 使得 $\alpha|_Z=\beta|_Z$。（在 \S11.3.3 中，我们将把它改进为：在良好的%
情形下——即 $Y$ \trSeparated{}时——只须 $\alpha|_{U\cap V}=\beta|_{U\cap V}$。）人们常把%
「映射」一词当作「态射」的同义词使用，这相当合理；只不过\trRationalMap{}未必是一个映射。因此为了%
避免混淆，当你的意思是「\trRationalMap{}」时，切不可只说「映射」。

我们也将讨论\trScheme{} $S$ 上诸 $S$-\trScheme{}的\trRationalMap{}。定义照旧，只是现在 $X$ 与 $Y$ %
都是 $S$-\trScheme，而 $\alpha\colon U\to Y$ 是 $S$-\trScheme{}之间的一个\trMorphism。

\textbf{7.5.2.}\ \hard\ \emph{更一般的\trRationalMap{}。}\quad 正如有理函数那样，定义~7.5.1 %
也可以推广到 $X$ 不\trReduced{}的情形，做法要么照搬（沿用同一个名字「\trRationalMap{}」），要么%
采用一个在非\trReduced\trLocus{}上施加某种控制的版本（\textbf{拟态射}（pseudo-morphisms），
[Stacks, tag 01RX]）。我们将在 \S11.3 中看到，从\trReduced{}\trScheme{}到\trSeparated{}%
\trScheme{}的\trRationalMap{}表现得格外好，这正是人们通常在这一语境下考虑它们的原因。拟态射这一%
定义的用意，便是把这些结果推广到 $X$ 非\trReduced{}的情形。我们不会用到拟态射这个观念。

\textbf{7.5.3.}\quad \trRationalMap{}的一个显然例子是\trMorphism{}本身。另一个重要的例子是投影%
$\mathbb{P}^n_A\dashrightarrow\mathbb{P}^{n-1}_A$，它由 $[x_0,\dots,x_n]\mapsto[x_0,\dots,x_{n-1}]$ 给出。%
（确切地说，按定义~7.5.1 的意义，这凭什么是一个\trRationalMap{}？）

\textbf{7.5.A.（练习）}\quad 设 $X$ 是一个\trReduced{}\trScheme。请描述 $X\dashrightarrow%
\mathbb{A}^1$ 型的\trRationalMap{}应当如何与 $X$ 上的\trRationalFunction{}相认同。

\trRationalMap{} $\pi\colon X\dashrightarrow Y$ 称为\textbf{\trDominant{}}（dominant）的（在某些文献里%
亦作 dominating），若对其某一个代表 $U\to Y$——从而也就是对其每一个代表——其像在 $Y$ 中%
\trDense。一个\trMorphism{}称为\textbf{\trDominant{}态射}（dominant morphism，或 dominating %
morphism），若它作为\trRationalMap{}是\trDominant{}的。

\textbf{7.5.B.（练习）}\quad 证明：\trIntegral{}\trScheme{}之间的\trRationalMap{}%
$\pi\colon X\dashrightarrow Y$ 是\trDominant{}的，当且仅当 $\pi$ 把 $X$ 的\trGenericPoint{}送到%
$Y$ 的\trGenericPoint。

\textbf{7.5.4.}\ \emph{\trRationalMap{}的复合。}\quad 稍加思索你便会相信：可以把一个从\trReduced{}%
\trScheme{} $X$ 到\trReduced{}\trScheme{} $Y$ 的\trDominant{}映射 $\pi\colon X\dashrightarrow Y$，%
与一个\trRationalMap{} $\rho\colon Y\dashrightarrow Z$ 复合起来（并且复合是\trWellDefined{}的）。%
此外，若 $\rho$ 是\trDominant{}的，则复合 $\rho\circ\pi$ 也是\trDominant{}的。\trIntegral{}\trScheme{}连同%
其间的\trDominant{}\trRationalMap{}构成一个在几何上饶有趣味的\trCategory。

\textbf{7.5.C.（容易的练习）}\quad 证明：\trIntegral{}\trScheme{}之间的\trDominant{}\trRationalMap{}按%
相反的方向给出\trFunctionField{}之间的\trMorphism。

在「适当经典的情形」（\trIntegral\trFiniteType{}的 $k$-\trScheme——特别地，还有将在%
\S11.2.9 中定义的\trIrreducible\trVariety）中，这件事是可逆的：\trDominant{}\trRationalMap{}对应于%
反方向的函数域的包含。我们将在下文 \S7.5.7 中把这一点说精确。但是，「函数域之间的%
\trMorphism{}总能给出\trDominant{}\trRationalMap{}」并不成立，甚至未必能给出\trRationalMap{}。例如，$\Spec k[x]$ %
与 $\Spec k(x)$ 具有同一个\trFunctionField{} $k(x)$，却不存在与之相应的 $k$-\trScheme{}%
的\trRationalMap{} $\Spec k[x]\dashrightarrow\Spec k(x)$。理由：这样的\trRationalMap{}应当对应于一个从%
$\Spec k[x]$ 的某个\trOpenSubset{} $U$（比如说 $\Spec k[x,1/f(x)]$）到 $\Spec k(x)$ %
的\trMorphism。但对任何一个 $f(x)$ 而言，（把 $k$ 逐点恒同地送到 $k$、把 $x$ 送到 $x$ %
的）环映射 $k(x)\to k[x,1/f(x)]$ 都不存在。

（如果你还想要更多的证据，来说明这个用拓扑方式定义的\trDominant{}观念同时也是代数的，你可以%
证明：若 $\phi\colon A\to B$ 是一个\trRingMorphism，则相应的\trMorphism{}%
$\Spec B\to\Spec A$ 是\trDominant{}的，当且仅当 $\phi$ 的\trKernel{}含于 $A$ 的\trNilradical{}%
之中。）

\textbf{7.5.5.}\ \emph{\trhDefinition{}。}\quad \trIntegral{}\trScheme{}之间的\trRationalMap{}%
$\pi\colon X\dashrightarrow Y$ 称为\textbf{\trBirational{}}（birational）的，若它是\trDominant{}的，并且%
存在另一个同样是\trDominant{}的\trRationalMap{}（一个「有理逆」）$\psi$，使得 $\pi\circ\psi$（在同一个%
\trEquivalenceClass{}中）就是 $Y$ 上的恒同，且 $\psi\circ\pi$（在同一个\trEquivalenceClass{}%
中）就是 $X$ 上的恒同。这就是「\trIntegral{}\trScheme{}连同其间的\trDominant{}\trRationalMap{}」这一%
\trCategory{}中的\trIsomorphism{}观念。（在可微范畴里，这并不特别有趣。）一个\trMorphism{}%
若作为\trRationalMap{}是\trBirational{}的，便称为\textbf{\trBirational{}态射}（birational morphism）。（注意：%
\trBirational{}态射的「逆」可以只是一个\trRationalMap{}，而不是\trMorphism。）

若存在一个双\trRationalMap{} $X\dashrightarrow Y$，我们便说 $X$ 与 $Y$（相互）\trBirational{}。%
若 $X$ 与 $Y$ 都\trIrreducible，则双\trRationalMap{}诱导\trFunctionField{}之间的\trIsomorphism。%
「对\trIntegral\trFiniteType{}的 $k$-\trScheme，函数域之间的映射对应于反方向的\trRationalMap{}」%
这一事实（将于命题~7.5.8 中证明）表明：在\trIntegral\trFiniteType{}的 $k$-\trScheme{}%
之间，凡诱导\trFunctionField{}的\trIsomorphism{}的映射必是\trBirational{}的。一个%
\trIntegral\trFiniteType{}的 $k$-\trScheme{}若是\trBirational{}于某个 $\mathbb{A}^n_k$，便称它是%
\textbf{有理}（rational）的。

\textbf{7.5.6.\ \trhProposition。}\ ——\ \emph{设 $X$ 与 $Y$ 是既\trIrreducible{}又%
\trReduced{}的\trScheme（据练习~5.2.G，即\trIntegral{}的概形）。则 $X$ 与 $Y$ \trBirational{}，%
当且仅当存在 $X$ 的一个\trDense\trOpenSubscheme{} $U$ 与 $Y$ 的一个\trDense%
\trOpenSubscheme{} $V$，使得 $U\cong V$。}

（我们加上「整」的假设只是为了简单。配以足够的细心与修改过的定义，这些假设是可以去掉的。）

命题~7.5.6 告诉你该如何看待双\trRationalMap{}。正如\trRationalMap{}是「大半处有定义的函数」，两个\trBirational{}%
的\trReduced{}\trScheme{}便是「大半处同构」的。例如，一个\trReduced\trFiniteType{}的%
$k$-\trScheme（比如一个\trAffine{} $k$-\trVariety）是\trRational{}的，当且仅当它有一个%
\trIsomorphic{}于 $\mathbb{A}^n_k$ 的某个\trOpenSubscheme{}的\trDense\trOpenSubscheme。

\emph{证明。}\ 「当」方向是显然的，故我们证明「仅当」方向。我们基本上就是顺着鼻子走，%
把手头给的条件全部用上（见图~\ref{fig:birational-proof-outline}）。

\begin{figure}[htbp]
\centering
{\scriptsize
\begin{tikzcd}[column sep=1.35em, row sep=2.55em]
f & g\circ f & g\circ f=\id\ \text{作为：} & & & g & f\circ g & f\circ g=\id\ \text{作为：} \\
\text{\trRationalMap{}} & \text{\trRationalMap{}} & \text{\trRationalMap{}} &
X \ar[r,<->,dashed,"\sim"] & Y &
\text{\trRationalMap{}} & \text{\trRationalMap{}} & \text{\trRationalMap{}} \\
\text{态射} & \text{”} & \text{”} &
X_1 \ar[u,hook'] \ar[ur] & Y_1 \ar[u,hook'] \ar[ul] &
\text{态射} & \text{”} & \text{”} \\
\text{”} & \text{态射} & \text{”} &
X_2 \ar[u,hook'] \ar[ur] & & & & \\
\text{”} & \text{”} & \text{态射} &
X_3 \ar[u,hook'] & Y_2 \ar[uu,hook'] \ar[l] &
\text{”} & \text{态射} & \text{”} \\
\text{”} & \text{”} & \text{态射} &
X_4 \ar[u,hook'] \ar[r,<->,"\sim"] & Y_3 \ar[u,hook'] &
\text{”} & \text{”} & \text{态射}
\end{tikzcd}}
\caption{命题~7.5.6 的证明概观。（依原书手绘图重绘。）}
\label{fig:birational-proof-outline}
\end{figure}

设 $F\colon X\dashrightarrow Y$ 与 $G\colon Y\dashrightarrow X$ 是一对互逆的\trRationalMap{}。%
选取\trRationalMap{} $F$ 与 $G$ 各自的代表\trMorphism{} $f\colon X_1\to Y$（其中 $X_1\subset X$）%
与 $g\colon Y_1\to X$（其中 $Y_1\subset Y$）。令 $X_2=(f|_{X_1})^{-1}(Y_1)\subset X_1$，%
则 $g\circ f$ 是一个从 $X_2$ 到 $X$ 的\trMorphism，它作为\trRationalMap{}是恒同。于是存在一个%
\trDense\trOpenSubset{} $X_3\subset X_2$，使得\trMorphism{} $g\circ f\colon X_3\to X$ %
是到 $X_3$ 上的恒同\trMorphism（或者更确切地说，它是\trOpenEmbedding{}%
$X_3\hookrightarrow X$）。

类似地，令 $Y_2=(g|_{Y_1})^{-1}(X_3)\subset Y_1$，于是 $f\circ g$ 是一个从 $Y_2$ 到 $Y$ %
的\trMorphism，它作为\trRationalMap{}是恒同。再令 $Y_3\subset Y_2$ 为一个\trDense\trOpenSubset，%
使得 $f\circ g\colon Y_3\to Y$ 是包含映射（即「恒同」）。

最后，令 $X_4=(f|_{X_3})^{-1}(Y_3)\subset X_3$，则 $(g\circ f)|_{X_4}$（经由 $Y_3$）%
是 $X_4$ 上的恒同\trMorphism，而 $(f\circ g)|_{Y_3}$（经由 $X_4$）是 $Y_3$ 上的恒同%
\trMorphism，于是我们便找到了这个开集之间的\trIsomorphism，它正是我们这个双\trRationalMap{}的%
一个「代表」。\hfill$\square$

\textbf{7.5.7.\ \trIrreducible\trVariety{}的\trRationalMap{}。}

\textbf{7.5.8.\ \trhProposition。}\ ——\ \emph{设 $X$ 是一个\trIntegral{}的 $k$-\trScheme，%
$Y$ 是一个\trIntegral\trFiniteType{}的 $k$-\trScheme，并给定一个保持 $k$ 的%
\trFunctionField{}扩张 $\phi^{\sharp}\colon K(Y)\hookrightarrow K(X)$。则存在 $k$-\trScheme{}%
的一个\trDominant{}\trRationalMap{} $\phi\colon X\dashrightarrow Y$ 诱导 $\phi^{\sharp}$。}

\emph{证明。}\ 用一个\trOpenSubset{}替换 $Y$ 之后，不妨设 $Y$ 是\trAffine{}的，比如说%
$\Spec B$，其中 $B$ 由有限多个元素 $y_1,\dots,y_n$ 在 $k$ 上生成。由于我们只须在 $X$ %
的一个\trOpenSubset{}上定义 $\phi$，同样不妨设 $X=\Spec A$ 是\trAffine{}的。于是%
$\phi^{\sharp}$ 给出一个包含 $\phi^{\sharp}\colon B\hookrightarrow K(A)$。把 $\phi^{\sharp}(y_i)$ %
写成 $f_i/g_i$（其中 $f_i,g_i\in A$），并令 $g:=\prod_i g_i$。于是 $\phi^{\sharp}$ 进而%
诱导一个包含 $B\hookrightarrow A_g$。从而 $\phi\colon\Spec A_g\to\Spec B$ 诱导%
$\phi^{\sharp}$。\trMorphism{} $\phi$ 是\trDominant{}的，因为零\trIdeal{}在包含 $B\hookrightarrow A_g$ %
之下的逆像仍是零\trIdeal，于是 $\phi$ 把 $X$ 的\trGenericPoint{}送到 $Y$ 的%
\trGenericPoint。\hfill$\square$

\textbf{7.5.D.（练习）}\quad 设 $K$ 是 $k$ 的一个\trFinitelyGenerated\trFieldExtension。%
（回想：$k$ 上的域扩张 $K$ 称为\textbf{有限生成}的，若 $K$ 中存在一个有限的「生成集」%
$\{x_1,\dots,x_n\}$，使得 $K$ 的每个元素都能写成系数属于 $k$ 的、以 $x_1,\dots,x_n$ 为%
变元的\trRationalFunction。）证明：存在以 $K$ 为\trFunctionField{}的\trIrreducible%
\trAffine{} $k$-\trVariety。（提示：考虑由 $t_i\mapsto x_i$ 给出的映射%
$k[t_1,\dots,t_n]\to K$，证明其\trKernel{}是一个\trPrimeIdeal{} $\mathfrak{p}$，且%
$k[t_1,\dots,t_n]/\mathfrak{p}$ 的\trFractionField{}为 $K$。几何解读：考虑由环映射%
$t_i\mapsto x_i$ 给出的映射 $\Spec K\to\Spec k[t_1,\dots,t_n]$，并取这个单点像的%
\trClosure。）

\textbf{7.5.E.（练习）}\quad 请描述下列诸范畴之间的\trEquivalenceOfCategories。

(a)~以「\trIntegral\trAffine{} $k$-\trVariety」为\trObject、「在 $k$ 上定义的\trDominant{}有理%
映射」为\trMorphism{}的范畴；以及

(b)~以「$k$ 的\trFinitelyGenerated\trFieldExtension」为\trObject、「在 $k$ 上延拓恒同%
的包含」为\trMorphism{}的（「箭头掉头」的）\trOppositeCategory。

（待我们在 \S11.2.9 中定义了一般的\trVariety{}之后，还可以添上：(c)~以「\trIntegral{} %
$k$-\trVariety」为\trObject、「在 $k$ 上定义的\trDominant{}\trRationalMap{}」为\trMorphism{}的范畴。）

特别地，一个\trIntegral\trAffine{} $k$-\trVariety{} $X$ 是\trRational{}的，若其\trFunctionField{}%
$K(X)$ 是 $k$ 的一个\textbf{纯超越扩张}（purely transcendental extension），也就是说%
$K(X)\cong k(x_1,\dots,x_n)$（对某个 $n$）。（这还需要说得更精确些：由 $X\to\Spec k$ %
诱导的映射 $k\leftarrow K(X)$，应当在这个\trIsomorphism{}之下与「显然的」映射%
$k\leftarrow k(x_1,\dots,x_n)$ 相合。）

\textbf{7.5.9.\ \trRationalMap{}的更多例子。}

这些例子里反复出现的一个主题是：到\trProjectiveScheme{}的\trRationalMap{}，其\trDomainOfDefinition{}%
可以越过余维数为 $1$ 的\trRegular{}点向上延拓。等我们讨论\trCurve{}时，将在「曲线到射影%
扩张定理」（Curve-to-Projective Extension Theorem）~15.3.1 中把这一点说精确。

\begin{figure}[htbp]
\centering
\begin{tikzpicture}[scale=1.15]
  \draw[thick] (0,0) circle (1);
  \draw[->] (-1.7,0) -- (1.85,0) node[right] {$x$};
  \draw[->] (0,-1.7) -- (0,1.85) node[right] {$y$};
  \fill (1,0) circle (1.7pt) node[above right] {$p$};
  \fill (-0.45,0.89) circle (1.7pt) node[above left] {$q$};
  \draw[thick] (-1.04,1.25) -- (1.65,-0.4);
  \node at (-1.2,-0.8) {$C$};
  \node at (1.5,-0.68) {斜率 $m$};
\end{tikzpicture}
\caption{用几何寻找本原勾股三元组。（依原书手绘图重绘。）}
\label{fig:pythagorean-stereographic}
\end{figure}

第一个例子是\textbf{勾股三元组}（Pythagorean triples）的经典公式，以及它经由「球极投影」%
（stereographic projection）的推导。设你正在寻找由 $x^2+y^2=1$ 给出的圆 $C$ %
（图~\ref{fig:pythagorean-stereographic}）上的\trRationalPoint。一个\trRationalPoint{}%
是 $p=(1,0)$。若 $q$ 是另一个有理点，则 $pq$ 是一条斜率为有理数（非无穷）的直线。这给出%
一个从\trConic{} $C$（现在解读为 $\Spec\mathbb{Q}[x,y]/(x^2+y^2-1)$）到 $\mathbb{A}^1_{\mathbb{Q}}$ %
的\trRationalMap{}，由 $(x,y)\mapsto y/(x-1)$ 给出。（刚才悄悄发生了一件微妙的事情：我们本来%
谈的是圆上的 $\mathbb{Q}$-点，结果却得到了概形之间的一个\trRationalMap{}。）反过来，给定过 $p$、%
斜率为 $m$（$m$ 为有理数）的一条直线，我们可以通过解方程组 $y=m(x-1)$ 与 $x^2+y^2=1$ %
把 $q$ 找回来。把第一个方程代入第二个，便得到一个关于 $x$ 的二次方程。我们知道 $x=1$ %
必是一个解（因为该直线与圆相交于 $(x,y)=(1,0)$），所以我们指望能把这个因子提出来，%
找出另一个因子。这确实行得通：
\[
\begin{gathered}
x^2+\bigl(m(x-1)\bigr)^2=1\\
\implies\quad (m^2+1)x^2+(-2m^2)x+(m^2-1)=0\\
\implies\quad (x-1)\bigl((m^2+1)x-(m^2-1)\bigr)=0.
\end{gathered}
\]
另一个解是 $x=(m^2-1)/(m^2+1)$，由此得 $y=-2m/(m^2+1)$。于是我们得到\trConic{} $C$ %
与以 $m$ 为坐标的 $\mathbb{A}^1$ 之间的一个双\trRationalMap{}，由 $f\colon(x,y)\mapsto y/(x-1)$ %
给出（它在 $x\neq 1$ 处有定义），其逆\trRationalMap{}由 $m\mapsto\bigl((m^2-1)/(m^2+1),\,-2m/(m^2+1)\bigr)$ %
给出（它在 $m^2+1=0$ 之外有定义）。

我们可以经由「包含」$\mathbb{A}^1_{\mathbb{Q}}\to\mathbb{P}^1_{\mathbb{Q}}$（日后我们将%
称之为一个\trOpenEmbedding）把它扩充为一个\trRationalMap{} $C\dashrightarrow\mathbb{P}^1_{\mathbb{Q}}$。%
此时 $f$ 由 $(x,y)\mapsto[y,\,x-1]$ 给出。于是我们有了一个有趣的问题：$f$ 的%
\trDomainOfDefinition{}是什么？看上去它处处有定义，唯独在 $y=x-1=0$ 之处——也就是除 $p$ %
之外——没有定义。但事实上它可以越过 $p$ 延拓！注意：$(x,y)\mapsto[x+1,\,-y]$（其中%
$(x,y)\neq(-1,0)$）在二者公共的\trDomainOfDefinition{}上与 $f$ 相合，因为%
$[x+1,\,-y]=[y,\,x-1]$。于是这个\trRationalMap{}能延拓得比我们起初所想的更远。这将是曲线到射影%
扩张定理~15.3.1 的一个特例。

\textbf{7.5.F.（练习）}\quad 利用上述办法，找出一个给出全体勾股三元组的「公式」。

\textbf{7.5.G.（练习）}\quad 证明：对任何\trCharacteristic{}不为 $2$ 的\trField{} $k$，%
$\mathbb{P}^2_k$ 中的\trConic{} $x^2+y^2=z^2$ 都\trIsomorphic{}于 $\mathbb{P}^1_k$。%
（如有需要，「$x^2+y^2=z^2$ 在 $\mathbb{P}^2_k$ 中」的定义可参\cref{ex:vanishing-scheme-definition}。）

\textbf{7.5.10.}\quad 事实上，$\mathbb{P}^2_k$ 中任何带有一个秩为 $3$ 的 $k$-\trValuedPoint{}%
（即\trResidueField{}为 $k$ 的一个点）的\trConic{}都\trIsomorphic{}于 $\mathbb{P}^1_k$ %
（要在基变换到 $\overline{k}$ 之后「秩」才有意义；参练习~5.4.J）。（带有 $k$-\trValuedPoint{}%
这一假设当然是必要的：$k=\mathbb{R}$ 上的 $x^2+y^2+z^2=0$ 便是一个不\trIsomorphic{}于%
$\mathbb{P}^1_k$ 的\trConic。）

\textbf{7.5.H.（练习）}\quad 通过找一个到 $\mathbb{A}^1_{\mathbb{Q}}$ 的双\trRationalMap{}——%
照搬处理\trConic{}时行之有效的那一套——求出 $y^2=x^3+x^2$ 的全部有理解。提示：你该从%
哪个点出发作投影？（在练习~19.10.I 中我们将会看到：这些点基本上构成一个\trGroup，而且%
这是一条退化的\trEllipticCurve。）

你将得到一个到 $\mathbb{P}^1_{\mathbb{Q}}$ 的\trRationalMap{}，它在「结点」（node）$x=y=0$ 之上%
没有定义，而且不能越过这个余维数 $1$ 的集合延拓。这恰好展示了我们日后的结果——讲述如何%
把\trRationalMap{}越过余维数 $1$ 的集合延拓到射影空间的曲线到射影扩张定理~15.3.1——的局限：%
余维数 $1$ 的集合必须是\trRegular{}的。

\textbf{7.5.I.（练习）}\quad 用类似的想法，找出从\trQuadricSurface{}%
$Q=\{x^2+y^2=w^2+z^2\}\subset\mathbb{P}^3_{\mathbb{Q}}$ 到 $\mathbb{P}^2_{\mathbb{Q}}$ %
的一个双\trRationalMap{}。再用它求出 $Q$ 上全部的\trRationalPoint。（这展示了解\textbf{丢番图%
方程}（Diophantine equations）的一种好办法。你将会找到 $Q$ 的一个\trIsomorphic{}于%
$\mathbb{P}^2$ 的某个\trDense\trOpenSubset{}的\trDense\trOpenSubset，在那里你可以%
轻轻松松求出全部\trRationalPoint。$Q$ 还会有一个\trClosedSubset，使得\trRationalMap{}在其上%
没有定义、或者不是\trIsomorphism，但你可以就事论事（ad hoc）地处置这个子集。）

\textbf{7.5.J.（练习：\trCremona{}（Cremona）变换——一个有用的经典构造）}\quad 考虑由%
$[x,y,z]\mapsto[1/x,\,1/y,\,1/z]$ 给出的\trRationalMap{} $\mathbb{P}^2_k\dashrightarrow\mathbb{P}^2_k$。%
它的\trDomainOfDefinition{}是什么？（它比 $xyz\neq 0$ 的\trLocus{}要大！）你将会观察到：%
你可以把它越过「余维数 $1$ 的集合」延拓（姑且不管我们迄今还不知道余维数是什么意思）。%
这又一次预示了曲线到射影扩张定理~15.3.1。

\textbf{7.5.11.}\ \hard\ \textbf{非\trRational{}的复\trCurve{}（有趣但非必要）。}\quad 我们现在来%
描述两个不存在从 $\mathbb{P}^1_{\mathbb{C}}$ 出发的非常值\trRationalMap{}的\trCurve{} $C$ 的例子。%
（诚然，我们迄今还不知道「\trCurve」是什么意思，但无妨。）两个证明用的都是\trFermat{}（Fermat）%
的\textbf{无穷递降}（infinite descent）法。这些结果可以（日后你将能借助定理~16.3.3 来%
核验）解读为这样一个事实：这些曲线没有「非平凡」的 $\mathbb{C}(t)$-\trValuedPoint——%
我们的意思是，任何 $\mathbb{C}(t)$-\trValuedPoint{}都暗中是一个 $\mathbb{C}$-\trValuedPoint。%
（你或许会注意到：若把同样的例子中的 $\mathbb{C}(t)$ 换成 $\mathbb{Q}$（并且 $C$ 是%
$\mathbb{Q}$ 而非 $\mathbb{C}$ 上的曲线），你就得到了数论与几何中的两个基本问题。%
练习~7.5.M 的类比是\trEllipticCurve{}上\trRationalPoint{}的问题，而你也许会意识到，%
练习~7.5.K 的类比甚至还要更出名。此外，练习~7.5.M(a) 在算术上的类比是「四平方数定理」%
（Four Squares Theorem）——不存在四个构成等差数列的整数平方——由\trFermat{}首次陈述。）

这些例子将让你得以一窥：数域上的事实为何、又如何常常有曲线的\trFunctionField{}上的%
事实与之平行。这种平行性是本学科中一个反复出现的深刻主题。

\textbf{7.5.K.（练习）}\quad 设 $n>2$，证明：从 $\mathbb{P}^1_{\mathbb{C}}$ 到%
$\mathbb{P}^2_{\mathbb{C}}$ 中的「\trFermat{}曲线」$x^n+y^n=z^n$ 不存在\trDominant{}\trRationalMap{}。提示：%
把它化归为证明：$f(t)^n+g(t)^n=h(t)^n$ 不存在「非常值」的解 $\bigl(f(t),g(t),h(t)\bigr)$，%
这里 $f(t)$、$g(t)$、$h(t)$ 是 $t$ 的\trRationalFunction（也就是 $\mathbb{C}(t)$ 的%
元素）。通过清去分母，再把它化归为证明：当 $f(t)$、$g(t)$、$h(t)$ 为互素多项式时，%
无非常值的解。为此，假设有解存在，并取其中正次数最低者加以考虑。然后利用 $\mathbb{C}[t]$ %
是\trUFD{}这一事实，以及等式 $h(t)^n-g(t)^n=\prod_{i=1}^{n}\bigl(h(t)-\zeta^i g(t)\bigr)$ %
（其中 $\zeta$ 是本原 $n$ 次单位根）。论证每个 $h(t)-\zeta^i g(t)$ 都是一个 $n$ 次幂。%
然后用
\[
\bigl(h(t)-g(t)\bigr)+\alpha\bigl(h(t)-\zeta g(t)\bigr)=\beta\bigl(h(t)-\zeta^2 g(t)\bigr)
\]
（其中 $\alpha$、$\beta$ 适当选取）来得到一个次数更低的解。（当 $n=2$ 时，这个论证%
为何失效？）

\textbf{7.5.L.（练习）}\quad 给出两条\trSmooth{}复\trCurve{} $X$ 与 $Y$，使得 $X$ 的任何%
非空\trOpenSubset{}都不\trIsomorphic{}于 $Y$ 的任何非空\trOpenSubset。（请试着不要为%
我们尚未定义「\trSmoothness」这件事而烦恼。）提示：练习~7.5.K。

\textbf{7.5.M.（练习）}\quad 设 $a$、$b$、$c$ 是互不相同的复数。请按下列步骤证明：若%
$x(t)$ 与 $y(t)$ 是 $t$ 的两个\trRationalFunction，使得
\begin{equation*}\tag{7.5.11.1}\label{eq:y-squared-cubic-product}
y(t)^2=\bigl(x(t)-a\bigr)\bigl(x(t)-b\bigr)\bigl(x(t)-c\bigr),
\end{equation*}
则 $x(t)$ 与 $y(t)$ 都是常数（即 $x(t),y(t)\in\mathbb{C}$）。（这里的 $\mathbb{C}$ 可以%
换成任何\trCharacteristic{}不为 $2$ 的\trField{} $K$；若 $K$ 不\trAlgClosed，则还需%
稍加小心。）

(a)~设 $P,Q\in\mathbb{C}[t]$ 是互素的多项式，使得它们有四个线性组合是完全平方，且这些%
完全平方中的任何两个都互非常数倍。证明：$P$ 与 $Q$ 都是常数（即 $P,Q\in\mathbb{C}$）。%
提示：通过给 $P$、$Q$ 改名，证明你可以假定这些完全平方恰是 $P$、$Q$、$P-Q$、$P-\lambda Q$ %
（对某个 $\lambda\in\mathbb{C}$）。定义 $u$ 与 $v$ 分别为 $P$ 与 $Q$ 的平方根。证明：%
$u-v$、$u+v$、$u-\sqrt{\lambda}\,v$、$u+\sqrt{\lambda}\,v$ 都是完全平方，且 $u$ 与 $v$ %
互素。若 $P$ 与 $Q$ 不都是常数，注意 $0<\max(\deg u,\deg v)<\max(\deg P,\deg Q)$。%
从一开始就假定 $P$ 与 $Q$ 被选作反例时已使 $\max(\deg P,\deg Q)$ 达到最小，由此导出%
矛盾。（题外话：三个互不相同的线性组合完全可能是完全平方。这类例子本质上对应于%
$\mathbb{C}(t)$ 中的本原勾股三元组——你能看出是怎么回事吗？）

(b)~设 $(x,y)=(p/q,\,r/s)$ 是 \eqref{eq:y-squared-cubic-product} 的一个解，其中%
$p,q,r,s\in\mathbb{C}[t]$，且 $p/q$ 与 $r/s$ 已写成最简形式。清去分母，证明%
$r^2q^3=s^2(p-aq)(p-bq)(p-cq)$。证明 $s^2\mid q^3$ 且 $q^3\mid s^2$，从而对某个%
$\delta\in\mathbb{C}$ 有 $s^2=\delta q^3$。由 $r^2=\delta(p-aq)(p-bq)(p-cq)$，证明%
$(p-aq)$、$(p-bq)$、$(p-cq)$ 都是完全平方。再证明 $q$ 也是完全平方，然后援引 (a)。

练习~21.4.H 给出了处理练习~7.5.K 与练习~7.5.M 的一种好得多的几何途径。
